\documentclass[11pt]{article}
\usepackage[margin=1in]{geometry}
\usepackage{amsmath,amssymb,mathtools}
\usepackage{fontspec,unicode-math}
\setmainfont{Latin Modern Roman}
\setsansfont{Latin Modern Sans}
\setmathfont{Latin Modern Math}
\usepackage{xurl,hyperref}
\hypersetup{colorlinks=true,urlcolor=blue,linkcolor=blue}
\newcommand{\sref}[2]{\hyperlink{source:#1:#2}{[S#2]}}
\newcommand{\cataloguescope}[1]{\paragraph{Recorded mathematical question.}#1}
\begin{document}
\section{L. Sadun — Problem}\label{statement}

\subsection{Definitions and mathematical statement}

\paragraph{Shared notation.}$\mathbb N=\{1,2,\ldots\}$, and $\mathbb Z,\mathbb Q,\mathbb R,\mathbb C$ denote the integers, rationals, reals and complex numbers. Any other conventions are specified locally.


Matching rules specify permitted local arrangements of finitely many labeled prototile types. For a type $t$, let $d(t)$ be its number density per unit volume in the resulting tilings. If $R$ is a finite union of tiles, write $N_t(R)$ for the number of its type-$t$ tiles and
\[D_t(R)=|N_t(R)-d(t)\operatorname{vol}(R)|.\]
Write $|\partial R|$ for $(d-1)$-dimensional boundary measure. Condition (A), low fluctuation, requires a constant $c_t$ for each tile type, uniform over finite regions $R$ in the tilings, such that $D_t(R)\leq c_t|\partial R|$. Condition (B) requires the same bound for every finite patch satisfying the local matching rules, including patches that need not extend to a complete admissible tiling. The densities and constants belong to the chosen rule system, rather than varying with $R$.
\paragraph{Statement recorded in the supplied source.}
Find matching rules in dimension $2$ satisfying condition (A), without requiring condition (B). \sref{statement}{2}
\subsection{Short English statement}
Find planar matching rules giving a boundary-size discrepancy bound for finite regions of complete tilings.
\subsection{Sources}
\begin{description}
\item[\hypertarget{source:statement:1}{S1}] ULAM.AI and the UnsolvedMath Contributors, \emph{UnsolvedMath: A Curated Collection of Open Mathematics Problems}, catalogue entry. CC BY 4.0. \url{https://huggingface.co/datasets/ulamai/UnsolvedMath}.
\item[\hypertarget{source:statement:2}{S2}] Faustin Adiceam (compiler), \emph{Open Problems and Conjectures related to the Theory of Mathematical Quasicrystals}, arXiv:1604.06280v2 (2016), Section 2.7.2, condition (A) and Problem 2.7.4, p. 10. \url{https://arxiv.org/pdf/1604.06280}.
\end{description}
\label{end:statement}
\end{document}
